\chapter{RSA Security and the Shared-Prime Failure}

\section{A Story About Ideal RSA vs. Real RSA}
In the ideal RSA story, each user generates two secret primes \(p\) and \(q\),
builds
\[
n=pq,
\]
and publishes \(n\) as part of the public key. Security then relies on one
assumption: given only \(n\), factoring it into \(p\) and \(q\) is hard.

If that assumption holds, the private key is safe. If it fails, RSA fails.

Now comes the real-world danger: two different public keys may accidentally
share a prime factor because of weak randomness during key generation. In that
case, the attack is no longer hard at all.

\section{How Shared Primes Break RSA}
Suppose two public moduli are
\[
n_1=pq_1,\qquad n_2=pq_2,
\]
with the same prime \(p\) reused by mistake.

Then an attacker computes
\[
\gcd(n_1,n_2)=p.
\]
Once \(p\) is known, both moduli are factored immediately:
\[
q_1=\frac{n_1}{p},\qquad q_2=\frac{n_2}{p}.
\]
After that, the private exponents can be recovered in the usual RSA way.

So the security collapse is dramatic: instead of solving a hard factoring
problem for each key, the attacker only computes a greatest common divisor.

\section{Why This Is Practical at Scale}

The key lesson is simple:
\begin{itemize}
  \item RSA mathematics is strong.
  \item RSA implementations can still fail if randomness is weak.
  \item Shared-prime detection with GCD is fast and devastating.
\end{itemize}

\section{How the Euclidean Algorithm Finds \(\gcd(n_1,n_2)\)}
\epigraph{The Euclidean algorithm is the granddaddy of all algorithms, because it is the oldest nontrivial algorithm that has survived to the present day.}{\textit{Donald Knuth}}

\index{Euclidean Algorithm}
Ironically, this same ``granddaddy'' algorithm is still a frontline tool more
than two thousand years later.

The 2012 study ``Ron was wrong, Whit is right'' analyzed large collections of
public keys and found many vulnerable keys due to shared factors, mostly caused
by insufficient entropy in embedded systems \cite{lenstra:ron-was-wrong}.

The Euclidean algorithm computes \(\gcd(a,b)\) efficiently by repeated
division-with-remainder:
\[
a=bq_0+r_0,\quad
b=r_0q_1+r_1,\quad
r_0=r_1q_2+r_2,\ \dots
\]
until the remainder becomes \(0\). The last nonzero remainder is the GCD.

\subsection*{Small worked example}
Take
\[
a=1071,\qquad b=462.
\]
Then
\[
1071=462\cdot2+147,
\]
\[
462=147\cdot3+21,
\]
\[
147=21\cdot7+0.
\]
Therefore
\[
\gcd(1071,462)=21.
\]

This is exactly the same operation used in the RSA shared-prime attack:
replace \(a,b\) by \(n_1,n_2\), and if the result is greater than 1, the two
keys share a factor and are broken.


\section*{References}
This chapter follows the large-scale RSA key study \cite{lenstra:ron-was-wrong}
and the Euclidean algorithm description in \cite{wiki:euclidean-algorithm}.
